\RequirePackage[T1]{fontenc}
\documentclass[conference,a4paper]{IEEEtran}
\usepackage{amsmath,amssymb,graphicx,booktabs,cite}
\usepackage[hidelinks]{hyperref}
\begin{document}
\typeout{SMOKE DIMENSIONS: paperwidth=\the\paperwidth;paperheight=\the\paperheight;textwidth=\the\textwidth;textheight=\the\textheight;columnsep=\the\columnsep;oddsidemargin=\the\oddsidemargin;topmargin=\the\topmargin;headheight=\the\headheight;headsep=\the\headsep}
\title{Manuscript Toolchain Verification}
\author{\IEEEauthorblockN{Isolated Test Document}\IEEEauthorblockA{Not a research manuscript}}
\maketitle
\begin{abstract}This isolated document verifies two-column rendering, mathematical glyphs and an embedded table. It contains no research findings.\end{abstract}
\begin{IEEEkeywords}Document preparation, verification\end{IEEEkeywords}
\section{Rendering Check}
An editable manuscript needs consistent text and equations. The capability variable is $\Delta e$, and the illustrative expression is
\begin{equation} e_i(t)=e_0+\Delta e.\end{equation}
\begin{table}[h]\centering\caption{Isolated check}\begin{tabular}{lr}\toprule Item&Value\\\midrule Columns&2\\Fonts&Embedded\\\bottomrule\end{tabular}\end{table}
\section{Interpretation}
Passing this check establishes only that the local renderer can produce a PDF from a small IEEE-style source. Final page count, content agreement and visual quality must still be checked on the finished manuscript.
\end{document}
